\documentclass[letterpaper]{article}
\usepackage[final]{rl_project}

\usepackage[utf8]{inputenc}
\usepackage[T1]{fontenc}
\usepackage{hyperref}
\usepackage{url}
\usepackage{booktabs}
\usepackage{amsfonts}
\usepackage{nicefrac}
\usepackage{microtype}
\usepackage{graphicx}

% Give your project report an appropriate title.
\title{Your Project Title}

% Replace the example names and email addresses with those of all group members.
% Authors sharing an affiliation can be listed together, as below.
\author{
  Joshua~B.~Evans \qquad Sophia~Jones \\
  Department of Computer Science \\
  University of Bath \\
  Bath, BA2 7AY \\
  \texttt{\{jbe25|sj2073\}@bath.ac.uk}
}

\begin{document}

\maketitle

\section{Problem Definition}

\section{Background}

\section{Method}

\section{Results}

\section{Discussion}

\section{Future Work}

\section{Personal Experience}

\section*{Generative AI Use Declaration}
% State whether your group used GenAI tools and, if so, how and to what extent.

\section*{References}
\small
% Include a complete list of references for sources cited in your report.

\normalsize
\clearpage

% Appendices
\appendix

\section{First Appendix Title}
Include supporting experimental details and results here, if needed.

\section{Second Appendix Title}
A report can have multiple appendices. Use them to organise your supplementary material sensibly.

\end{document}
